\section{Method 3: Rectified Flow Models}

% --- Slide 1: Rectified Flow Fundamentals ---
\begin{frame}{Rectified Flow Fundamentals}
    \vspace*{-0.15cm}
    \begin{center}
        \includegraphics[width=\linewidth, keepaspectratio]{diagram_rectified_flow_process.png}
    \end{center}

    \small
    \begin{itemize}
        \item \textbf{Training.} The image and noise mix on one straight line at a random time.
        \item \textbf{Generation.} The model predicts the direction, and the sampler steps along it.
        \item \textbf{Our setup.} A released checkpoint, already step distilled, so four steps are enough.
    \end{itemize}
\end{frame}

% --- Slide 2: Architecture, FLUX.2-klein plus LoRA ---
\begin{frame}{Architecture: FLUX.2-klein and LoRA}
    \vspace*{-0.15cm}
    \begin{center}
        \includegraphics[width=\linewidth, keepaspectratio]{diagram_architecture_lora.png}
    \end{center}

    \small
    \begin{itemize}
        \item \textbf{Flow model.} A 4 billion parameter rectified flow transformer.
        \item \textbf{Layers.} Five double-stream blocks, then 20 single-stream blocks.
        \item \textbf{Adapter.} It adapts 80 weights, and the base weights stay frozen.
    \end{itemize}
\end{frame}

% --- Slide 3: The Generation Pipeline on 4 GB ---
\begin{frame}{The Generation Pipeline on 4 GB}
    \vspace*{-0.25cm}
    \begin{center}
        \includegraphics[width=\linewidth, keepaspectratio]{diagram_rectified_flow_pipeline.png}

        \vspace{-0.05cm}

        \includegraphics[width=\linewidth, keepaspectratio]{diagram_generation_internals.png}
    \end{center}

    \vspace{-0.1cm}
    \small
    \begin{itemize}
        \item \textbf{Why it fits.} The text encoder runs once on the CPU, then its memory is freed.
        \item \textbf{Safety gate.} A frame with more than one face is refused, and there is no bypass.
    \end{itemize}
\end{frame}

% --- Slide 4: Cleanup Stages and Results ---
\begin{frame}{Cleanup Stages and Results}
    \vspace*{-0.15cm}
    \begin{center}
        \includegraphics[width=\linewidth, keepaspectratio]{diagram_cleanup_stages.png}
    \end{center}

    \vspace{0.1cm}
    \begin{center}
        {\scriptsize
        \begin{tabular}{lrrrrr}
            \toprule
            \textbf{Stage} & \textbf{Face change} & \textbf{p95} & \textbf{Consecutive identity} & \textbf{Identity to portrait} & \textbf{Hair noise} \\
            \midrule
            Control & 6.657 & 12.3488 & 0.985 & $-0.036$ & 31.600 \\
            Raw swap & 7.451 & 17.3794 & 0.923 & 0.215 & 32.570 \\
            Restoration & 7.160 & 16.2097 & 0.927 & 0.215 & 32.072 \\
            Flow on raw & 6.666 & 15.6581 & 0.939 & 0.215 & 32.060 \\
            Restoration + flow & 6.630 & 15.2802 & 0.938 & 0.215 & 31.731 \\
            Deflicker & 5.627 & 12.6650 & 0.936 & 0.211 & 31.896 \\
            \bottomrule
        \end{tabular}}

        \vspace{0.08cm}
        {\scriptsize Settings: strength 0.6 for flow on raw, 0.3 for restoration + flow. \\
        Reading: these metrics are temporal change indicators, not proof; motion, lighting, texture change and flicker all enter the number, so read them together with the visual observation.}
    \end{center}

    \vspace{0.05cm}
    \small
    \begin{itemize}
        \item \textbf{Recommended.} Restoration, then optical flow at strength 0.3.
        \item \textbf{Rejected.} The deflicker is cheaper, but its low face change suggests lost motion.
    \end{itemize}
\end{frame}

% --- Slide 4b: Reassembly and the Boil ---
\begin{frame}{Reassembly and the Boil}
    \vspace*{-0.1cm}
    \begin{center}
        \begin{minipage}[b]{4.2cm}
            \centering
            \animategraphics[autoplay,loop,width=\linewidth]{10}{anim_raw/frame_}{000}{031}\\[0.08cm]
            {\scriptsize Raw assembly}
        \end{minipage}\hfill
        \begin{minipage}[b]{4.2cm}
            \centering
            \animategraphics[autoplay,loop,width=\linewidth]{10}{anim_flow/frame_}{000}{031}\\[0.08cm]
            {\scriptsize Flow on the raw frames}
        \end{minipage}\hfill
        \begin{minipage}[b]{4.2cm}
            \centering
            \animategraphics[autoplay,loop,width=\linewidth]{10}{anim_chain/frame_}{000}{031}\\[0.08cm]
            {\scriptsize Restoration + flow}
        \end{minipage}
    \end{center}

    \vspace{0.1cm}
    \small
    \begin{itemize}
        \item \textbf{Reassembly.} H.264 at 29.969 fps, with the original audio muxed back.
        \item \textbf{Watch the boil.} It is clear in the first clip and reduced in the last one. Some frames still jump.
    \end{itemize}

    \note{The embedded PDF animation plays only in a viewer that supports animate.sty. A local MP4 backup of the clip exists.}
\end{frame}

% --- Slide 5: Adapter Versions and Failure Analysis ---
\begin{frame}{Adapter Versions and Failure Analysis}
    \vspace*{-0.2cm}
    \begin{columns}[T]
        \begin{column}{0.46\textwidth}
            \centering
            {\footnotesize
            \begin{tabular}{lll}
                \toprule
                 & \textbf{v1} & \textbf{v1.1} \\
                \midrule
                Seconds per frame & 16.34 & about 199 \\
                Peak VRAM & 3.60 GiB & \textbf{4.68 GiB} \\
                Adapter rank & 128 & \textbf{512} \\
                96 frames & 26 min & about 5.3 h \\
                \bottomrule
            \end{tabular}}
        \end{column}

        \begin{column}{0.50\textwidth}
            \small
            \begin{itemize}
                \item \textbf{Measured.} Only v1 ran the metric script.
                \item \textbf{Memory.} v1.1 went past the card, 12x slower.
                \item \textbf{Judgement.} We preferred v1 by eye, so the clip kept v1.
                \item \textbf{Tail.} The worst frames still jump.
            \end{itemize}
        \end{column}
    \end{columns}

    \vspace{-0.1cm}
    \begin{center}
        \includegraphics[width=\linewidth, keepaspectratio]{faceswap_lora_compare_grid_0270.png}

        \vspace{-0.1cm}
        {\scriptsize Five columns: input body crop, v1, v1.1 with prompt A, v1.1 with prompt B, and the reference portrait.}
    \end{center}
\end{frame}
